\begin{definition}
Let $n\ge2$, $p>1$ and $\varepsilon,\gamma_0,\gamma_1\in\R$, and let
$A\colon\R^n\to\R^n$.  The predicate
$\mathrm{CitedDirichletW1p}(n,p,\varepsilon,\gamma_0,\gamma_1,A)$ is the cited
result (C2) of the Note for the formalization (monotone operator theory):

\emph{if $A$ satisfies (Q1) of the thesis--(Q2) of the thesis with parameters
$p,\varepsilon,\gamma_0,\gamma_1$ for some $\varepsilon\ge0$
(\noderef{EllipticStructureQ12}) and $A(0)=0$, then for every ball
$B=B_R(x_0)$ and every $b\in W^{1,p}(B)$ with weak gradient $g_b$
(\noderef{MemW1pOn}) there is a weak solution $v$ of
$\operatorname{div}A(\nabla v)=0$ in $B$ (\noderef{IsWeakSolutionOn}), with weak
gradient $g_v$, such that $v-b\in W^{1,p}_0(B)$ (\noderef{MemW1pZeroOn}).}

This is used as an explicit hypothesis in the passage from $\varepsilon>0$ to
$\varepsilon=0$ and is never introduced as an axiom.
\end{definition}
